% =====================================================================
% SECCIÓN DEL INTEGRANTE 3: Análisis de complejidad de avanzar() y retroceder()
% Pensado para la plantilla UNL (formato.tex): usa booktabs y float, que ya
% vienen cargados ahí. Se pega dentro de la sección "Análisis de complejidad".
% =====================================================================

\subsection{Análisis de complejidad de \texttt{avanzar()} y \texttt{retroceder()}}

Los métodos \texttt{avanzar()} y \texttt{retroceder()} mueven el cursor
\texttt{actual} de la \texttt{ListaDobleNavegacion} un nodo hacia adelante o
hacia atrás. Cada nodo (\texttt{NodoDoble}) guarda dos referencias: una al nodo
\texttt{siguiente} y otra al nodo \texttt{anterior}. Gracias a esto, el cursor
puede moverse en ambos sentidos sin recorrer la lista.

\subsubsection*{Complejidad temporal: $O(1)$}

Cada método ejecuta siempre la misma cantidad de pasos, sin importar cuántos
nodos ($n$) tenga la lista:

\begin{enumerate}[leftmargin=*]
  \item Una comprobación: que \texttt{actual} no sea \texttt{null} y que exista
        el nodo vecino (\texttt{actual.siguiente} o \texttt{actual.anterior}).
  \item Una reasignación de referencia: \texttt{this.actual = this.actual.siguiente}
        (o \texttt{.anterior}).
  \item Un retorno: el dato del nuevo nodo, o \texttt{null} si se llegó a un
        extremo de la lista.
\end{enumerate}

No hay ciclos (\texttt{for}, \texttt{while}) ni recursión. El número de
operaciones es constante, por lo que $T(n) = c$ y la complejidad temporal es
$O(1)$ tanto en el mejor como en el peor caso.

\subsubsection*{Complejidad espacial: $O(1)$}

Ninguno de los dos métodos crea nodos, arreglos ni copias de los datos. Solo
cambian a qué nodo apunta la referencia \texttt{actual}, que ya existía. La
memoria adicional que usan es constante, así que la complejidad espacial
auxiliar es $O(1)$. Como solo se mueve una referencia, el Garbage Collector no
recibe trabajo extra por navegar.

\subsubsection*{Comparación con una lista simplemente enlazada}

En una lista simple cada nodo solo conoce al \texttt{siguiente}. Avanzar sigue
siendo $O(1)$, pero para retroceder habría que recorrer la lista desde la
cabeza hasta el nodo anterior al cursor, lo que en el peor caso cuesta $O(n)$.
El enlace \texttt{anterior} de la lista doble evita ese recorrido, a cambio de
una referencia extra por nodo.

% Tabla con el estilo de la plantilla UNL (booktabs, letra pequeña y leyenda)
\begin{table}[H]
    \centering\small
    \begin{tabular}{@{}lcc@{}}
    \toprule
        \textbf{Operación} & \textbf{Lista simple} & \textbf{Lista doble} \\
    \midrule
        \texttt{avanzar()}    & $O(1)$ & $O(1)$ \\
        \texttt{retroceder()} & $O(n)$ & $O(1)$ \\
        Memoria por nodo      & 1 referencia & 2 referencias \\
    \bottomrule
    \end{tabular}
    \caption{Comparación de la navegación en lista simple y lista doble.}
    \label{tab:navegacion}
\end{table}

\subsubsection*{Verificación con pruebas unitarias}

Las pruebas 2 y 3 de \texttt{tests/navegacion.test.js} comprueban este
comportamiento (ver Tabla \ref{tab:cobertura}) con una lista de tres registros (101, 102 y 103). La prueba 2
avanza hasta la cola y verifica que en el límite superior se devuelva
\texttt{null}. La prueba 3 avanza hasta la cola, retrocede hasta la cabeza y
verifica que en el límite inferior también se devuelva \texttt{null}, sin que
el cursor se pierda ni se lance un error.
